\usepackage{xeCJK}
\usepackage{amsmath,amssymb,amsfonts,amsthm}
\usepackage{physics}
\usepackage{indentfirst}
\usepackage{mathrsfs}
\usepackage{bbm}
\usepackage{graphicx}
\usepackage{epstopdf}
\usepackage{hyperref}
\usepackage{simpler-wick}
\usepackage{caption}
\usepackage{float}
\usepackage{textcomp}
\usepackage{geometry}
\geometry{textheight=230mm,textwidth=165mm,footskip=20mm}
% \geometry{b5paper}
\usepackage{parskip}
\parindent=20pt


\renewcommand{\thefootnote}{\fnsymbol{footnote}} 

\usepackage{epigraph} % for quote
\makeatletter
  \let\@epipart\@endpart
  \renewcommand{\@endpart}{\thispagestyle{epigraph}\@epipart}
\makeatother        

\usepackage{feynmf}


\usepackage{natbib} % 设置样式为 authoryear
\bibliographystyle{plain} % 指定 .bib 文件（需加扩展名）
\renewcommand{\bibname}{参考文献}

% \usepackage[backend=biber, style=numeric]{biblatex}
% \addbibresource{参考.bib}

\usepackage{slashed}



% terminology
\usepackage[acronym]{glossaries} % [acronym] 选项允许我们定义缩写词

\makeglossaries % 初始化glossaries系统，必须添加

% 定义中文术语


\newglossaryentry{scal}{
  name={标量},
  description={Scalar}
}

\newglossaryentry{vec}{
  name={矢量},
  description={Vector}
}

\newglossaryentry{spinor}{
  name={旋量},
  description={Spinor}
}

\newglossaryentry{局域性}{
  name={局域性},
  description={Locality}
}

\newglossaryentry{Lag}{
  name={Lagrangian},
  description={拉格朗日量}
}

\newglossaryentry{LzB}{
  name={Lorentz加速变换},
  description={Lorentz Boost}
}
